\section{Can the weighting learn its way out?}
\label{sec:extensions}

Everything to this point audits one fixed weighting: two parameters,
$\alpha$ and $p$, chosen by grid search per dataset. The objective's own
future-work section points past them --- weight each detail band separately,
adapt over image position, deepen the decomposition, learn the schedule from
$\sigma$ instead of assuming $\sigma^{p}$, even learn the basis itself. If the
shortfall of \S\ref{sec:worth} is an artefact of one hand-chosen weighting, a
learned or finer-grained one should escape it. This section turns each of
those axes on separately, so that a negative answer attaches to the idea
rather than to one implementation.

\paragraph{Protocol.} Every extension holds its summed band weight at $1$ for
all $\sigma$, so each is read against the normalised variants of
\S\ref{sec:confound} --- \emph{not} against the published form, whose
gradient-magnitude schedule normalisation removes, and whose comparison would
otherwise confound the new axis with a learning-rate change. Five seeds unless
marked; pairing, correction and metrics as throughout.

\begin{table}[t]
\centering
\small
\caption{The method extensions, CIFAR-10 epoch \num{199}. Reference rows give
the comparisons each extension must beat: the published weighting, and the
budget-held-fixed variant that isolates the frequency component.
$\Delta$ is the paired-by-seed difference from MSE; $p$ is Holm-corrected
across the comparison family of Tables~\ref{tab:compare-fid}
and~\ref{tab:confound}. A superscript marks arms run at three seeds.}
\label{tab:extensions}
\begin{tabular}{llrr}
\toprule
Arm & Axis turned on & $\Delta$ FID & $p$ (Holm) \\
\midrule
\multicolumn{4}{l}{\emph{reference}} \\
Published weighting & $\alpha$, $p$ by grid search & $+0.44$ & $0.43$ \\
\quad budget held fixed & shape alone & $+1.34$ & $0.004$ \\
\midrule
Per-subband powers & steeper schedule for HH & $+2.61$ & $0.001$ \\
Spatially adaptive & weights vary over position & $+4.65$ & $0.0002$ \\
Learned basis & lifting filters, Haar init & $+6.13$\textsuperscript{3} & $0.015$ \\
Learned schedule & $\sigma$-conditioned, no $\alpha$/$p$ & $+6.66$ & $0.0001$ \\
Multi-level (3) & deeper decomposition & $+88.87$ & $0.021$ \\
GradNorm balancing & gradients set the weights & $+330.39$ & $0.0002$ \\
All axes combined & learned, spatial, lifting & $+12.53$\textsuperscript{3} & $0.004$ \\
\bottomrule
\end{tabular}
\end{table}

No route escapes. Every extension in Table~\ref{tab:extensions} is
significantly worse than the MSE baseline, and every one is worse than the
budget-held-fixed weighting it was meant to improve upon --- the nearest of
them by $+1.3$~FID. The ordering is the finding: each additional degree of
adaptive freedom costs more than it returns, which is what
Eq.~\ref{eq:parseval} leads one to expect. Sub-band weighting reallocates a
fixed error budget; \S\ref{sec:confound} showed that the pure-frequency
version of the reallocation is already significantly harmful; the extensions
grant that reallocation more freedom still, and it uses them.

Two arms did not degrade but diverged. GradNorm balancing ends at
$348.8 \pm 28.0$~FID, joining SNR-weighted (\S\ref{sec:replication}) in
outright training collapse: both schemes reallocate gradient mass according
to their own dynamics while training runs, and both destabilise what they
are reweighting. The three-level decomposition is unstable rather than
collapsed --- its across-seed spread is two orders of magnitude above any
tier-1 arm's --- which is its own warning: depth buys variance before it
buys anything else.

The decisive row is the learned schedule. A $\sigma$-conditioned uncertainty
weighting removes $\alpha$ and $p$ entirely and lets training find the
schedule --- exactly the remedy proposed for having to tune those parameters
per dataset. It lands at $+6.66$~FID ($p = 0.0001$). What gradient descent
finds in this family is worse than what grid search published, and
\S\ref{sec:confound} explains why: whatever freedom the optimiser is given
is freedom over the frequency axis, and the frequency axis is the one that
hurts.

These are extensions of one objective at one resolution, not independent
members of the family; \S\ref{sec:limitations} carries that scope forward.
What they do establish is that the shortfall of \S\ref{sec:worth} is not an
artefact of under-engineered weights. There is no setting of this mechanism,
fixed or learned, in which we found the correction paying for itself.
